\section{Conclusion}\label{sec:conclusion}

This working paper filed the first half of a pre-registered
experiment on the predictability of Supreme Court votes. The
frozen corpus exists, is hash-chained, and is public; the
statistical baselines are committed, and the strongest of them ---
one bit of revealed ideology per justice---sets a vote-level bar
of 63.7\% on the held-out window; and a trained challenger,
armed with everything the structured case file knows before the
decision, failed to clear the bar with three model families and a
clean ablation explaining why. The predictable part of a modern
justice's vote, to the resolution that the tested metadata can
see, is the
justice.

The second half is where the project's real question lives. The
same corpus---its 793 distinct opinion texts collected, normalized,
and audited---
will face the sealed fifty with a zero-shot reader, a retrieval
reader, and nine per-justice persona adapters trained on each
justice's own past words---each condition pre-registered, each
adapter audited for memorization, and each result, positive or
null, filed as a finding. The experiment was designed so that
either answer will be worth reading: a persona that beats the
zero-shot reader on unseen future cases would be, to our knowledge,
the first
sealed-holdout evidence that a justice-specific textual signal
carries vote information; a persona that does not would be, to our
knowledge, the cleanest
evidence yet that, beyond ideology, the written record keeps its
counsel. The bar is set, the seal is frozen, and the remaining
variable is the text.
